\documentclass[11pt]{article}

\usepackage[margin=1in]{geometry}
\usepackage{amsmath,amssymb,amsthm}
\usepackage{enumitem}
\usepackage[hidelinks]{hyperref}
\urlstyle{same}

\theoremstyle{plain}
\newtheorem{theorem}{Theorem}
\newtheorem{proposition}{Proposition}
\newtheorem{lemma}{Lemma}
\newtheorem{corollary}{Corollary}
\theoremstyle{definition}
\newtheorem{definition}{Definition}
\newtheorem{remark}{Remark}

\newcommand{\temp}{\operatorname{temp}}
\newcommand{\mean}{\operatorname{mean}}
\newcommand{\den}{\operatorname{den}}
\newcommand{\gr}{\operatorname{gr}}
\newcommand{\Games}{\mathcal G}
\newcommand{\Numbers}{\mathcal N}
\newcommand{\Inf}{\mathcal I}
\newcommand{\Q}{\mathbb Q}

\title{The Thermic Degree of a Norton Product}
\author{a9lim}
\date{August 2026}

\begin{document}
\maketitle

\begin{abstract}
Let $G.u$ denote Norton multiplication of a short game by a positive dyadic
number $u=m/2^k$ in canonical form.  Put
$\delta_u=1/\den(u)=2^{-k}$ and $a_u=u-\delta_u$.  We prove the exact
thermic law
\[
  \mean(G.u)=u\mean(G),\qquad
  \temp(G.u)=u\temp(G)+a_u
\]
for every nonnumeric game in the ordinary finite thermographic domain.  For a
number $x$, whose canonical mesh is $\epsilon_x=1/\den(x)$, the second formula
is replaced by
\[
 \temp(x.u)=
 \begin{cases}
 -1,&a_u-u\epsilon_x<0,\\
 a_u-u\epsilon_x,&a_u-u\epsilon_x\geq0.
 \end{cases}
\]
Thus every positive numeric Norton operator descends to an additive transport
from thermic degree $\tau$ to degree $u\tau+a_u$.

These transports do not form a dyadic action.  Their exact degree defect is
\[
 \Delta(u,v)=v(1-\delta_u)-\delta_v+\delta_{uv}\geq0,
\]
and we classify all pairs for which it vanishes.  The formula, its numeric
branch, and the complete defect classification are the main results.  We end
by locating them inside the more modest tropical comparison: temperature is a
filtration degree and thermographic option folds use max-plus and min-plus,
but the freeze step is not tropical.  Moreover the degree-zero residue is the
group of infinitesimals and contains $[*]\neq0$ with $2[*]=0$, so no faithful
lift-compatible action of full dyadic coefficients can turn the filtration
into a Newton-style graded ring.
\end{abstract}

\section{Norton and thermographic conventions}

We work with short normal-play partizan games.  Numbers are assigned
temperature $-1$; every game discussed below is in the ordinary finite
thermographic domain.  The recursive definition and the standard order and
additivity properties of Norton multiplication may be found in
\cite[Ch.~8]{berlekamp2001}, \cite[\S6.2]{johnson2011}, and
\cite[Prop.~13]{mesdal2009}.

There is a form convention to make explicit.  Norton multiplication is
independent of the chosen form of its first argument, but its definition uses
the selected Left and Right options of the positive unit
\cite{moews2002}.  Different presentations of the same unit can therefore
give different recurrences: in the usual notation, for example,
$(1/2).\{1*\mid\}=\{1\mid0\}$ whereas
$(1/2).\{0\mid\}=1/2$.  Throughout this paper a numeric unit means its
\emph{canonical short-game form}.  This is also the convention of the shipped
implementation, which canonicalizes the unit before extracting its
incentives.

Let $u=m/2^k>0$ be in lowest terms, and define
\begin{equation}\label{eq:mesh-shift}
 \delta_u=\frac1{\den(u)}=2^{-k},
 \qquad a_u=u-\delta_u.
\end{equation}
For $k=0$ this says $\delta_u=1$.  The canonical options of a nonintegral $u$
are $u-\delta_u$ and $u+\delta_u$; a positive integer has the sole Left
option $u-1$.  Every Norton increment is consequently the same number $a_u$:
a Left option contributes $u-\delta_u$, and a Right option contributes
$2u-(u+\delta_u)$.  Writing $A_u(G)=G.u$, the defining recurrence reduces to
\begin{equation}\label{eq:norton-recursion}
 A_u(n)=nu\quad(n\in\mathbb Z),
 \qquad
 A_u(G)=\{A_u(G^L)+a_u\mid A_u(G^R)-a_u\}
 \quad(G\notin\mathbb Z).
\end{equation}

We next fix the thermographic notation used in the proof.  For a number $x$,
set $LS_t(x)=RS_t(x)=x$.  For a nonnumber $G$, form the raw option envelopes
\begin{equation}\label{eq:raw-walls}
 L_G^{\rm raw}(t)=\max_{G^L}RS_t(G^L),
 \qquad
 R_G^{\rm raw}(t)=\min_{G^R}LS_t(G^R).
\end{equation}
Before freezing, the left and right scaffold walls are
\begin{equation}\label{eq:scaffold-walls}
 LS_t^{\rm sc}(G)=L_G^{\rm raw}(t)-t,
 \qquad
 RS_t^{\rm sc}(G)=R_G^{\rm raw}(t)+t.
\end{equation}
They are followed until their first meeting and are then replaced by the
common constant mast.  The resulting frozen walls are denoted $LS_t(G)$ and
$RS_t(G)$; at $t=0$ they give the ordinary stops $LS(G)$ and $RS(G)$.  Thus
\begin{equation}\label{eq:gap}
 D_G(t)=L_G^{\rm raw}(t)-R_G^{\rm raw}(t)-2t
       =LS_t^{\rm sc}(G)-RS_t^{\rm sc}(G),
\end{equation}
and $\temp(G)$ is the least $t\geq0$ at which this raw gap closes.  Wall slopes
give that $D_G$ is nonincreasing, $LS_t$ is nonincreasing, and $RS_t$ is
nondecreasing.

For later use, if $X$ is a Left option in \eqref{eq:norton-recursion}, define
its \emph{transformed right wall at height $s$} by
\begin{equation}\label{eq:transformed-walls}
 \mathcal T^R_X(s)=RS_s(A_u(X))+a_u-s.
\end{equation}
This is exactly the candidate that $X$ contributes to the left scaffold of
$A_u(G)$.  Dually, a Right option $Y$ contributes the transformed left wall
\[
 \mathcal T^L_Y(s)=LS_s(A_u(Y))-a_u+s.
\]
The terminology will never refer merely to $RS_s(A_u(X))$ or
$LS_s(A_u(Y))$; the parent shift in \eqref{eq:transformed-walls} is part of the
definition.

We use number avoidance in its standard form: in a sum $H+x$ with $x$ a
number and $H$ a nonnumber, any player with an optimal move has one in the
$H$ component.  In particular, comparisons with a number can be proved by
giving the option-wise strategy in the nonnumeric component
\cite[Part~II]{siegel2013}.  We will also use the standard stop implication
\begin{equation}\label{eq:stop-implication}
 RS(H)<m\Longrightarrow H\not\geq m,
\end{equation}
which says equivalently that Right wins $H-m$ when moving first; the mirror
statement uses $LS(H)>m$.

\section{The exact positive-dyadic law}

The thermic formula appears not to have been located in the literature.
\footnote{The electronically accessible accounts in
\cite[\S6.2]{johnson2011} and \cite[Prop.~13]{mesdal2009} give Norton
multiplication and its standard algebraic properties, but not this affine
thermic law.  Two print-only loci remain to be checked by hand before any
stronger novelty claim is made: \emph{Winning Ways}, Chapter~8, and
Siegel, \S III.3.  Accordingly, ``not located'' is meant literally, not as a
claim of absolute novelty.}
The proof needs a separate numeric branch because Norton recursion stops at
integers, not at arbitrary numbers.

\begin{lemma}[numeric images]\label{lem:numeric-images}
Let $x$ be a short-game number and put
$\epsilon_x=1/\den(x)$.  Then
\[
 \mean(A_u(x))=ux,
 \qquad
 \temp(A_u(x))=
 \begin{cases}
 -1,&a_u-u\epsilon_x<0,\\
 a_u-u\epsilon_x,&a_u-u\epsilon_x\geq0.
 \end{cases}
\]
In particular $\temp(A_u(x))<a_u$.
\end{lemma}

\begin{proof}
Induct on the exponent of the exact dyadic denominator of $x$.  Integers are
the base clause of \eqref{eq:norton-recursion}, so their images are the numbers
$ux$ and have temperature $-1$.  If $x$ is nonintegral with canonical mesh
$\epsilon=\epsilon_x$, then
\[
 x=\{x-\epsilon\mid x+\epsilon\}.
\]
The two option meshes are $2\epsilon$ (unless an option is already integral),
so their candidate temperatures are
$a_u-2u\epsilon<a_u-u\epsilon$.  By induction the option images have masts
$u(x-\epsilon)$ and $u(x+\epsilon)$ and freeze strictly before the height
$c=a_u-u\epsilon$ whenever $c\geq0$.  At height $c$ the two parent scaffold
values are therefore
\[
 u(x-\epsilon)+a_u-c=ux,
 \qquad
 u(x+\epsilon)-a_u+c=ux.
\]
On the terminal interval immediately below $c$ their gap is $2(c-s)>0$.
Since the raw gap is nonincreasing, it could not have closed earlier.  Hence
the first meeting is at $c$ and its mast is $ux$.

If $c<0$, every lower-denominator option image is a number.  The two numeric
bounds in \eqref{eq:norton-recursion} are
$ux+c<ux<ux-c$; the simplicity rule selects the number $ux$.  This is the
temperature-$-1$ branch.  In both cases the displayed temperature is strictly
below $a_u$, since $u\epsilon>0$.
\end{proof}

The next identity is the algebraic center of the proof.  It is stated for the
raw gap because that is where the two parent shifts cancel exactly.

\begin{lemma}[affine raw-gap transport]\label{lem:raw-gap}
Let $G$ be nonintegral, and suppose that for every option $X$ of $G$ and every
$t\geq0$,
\[
 LS_{a_u+ut}(A_u(X))=uLS_t(X),
 \qquad
 RS_{a_u+ut}(A_u(X))=uRS_t(X).
\]
Then
\begin{align}
 L_{A_u(G)}^{\rm raw}(a_u+ut)
   &=uL_G^{\rm raw}(t)+a_u,\label{eq:left-raw-transform}\\
 R_{A_u(G)}^{\rm raw}(a_u+ut)
   &=uR_G^{\rm raw}(t)-a_u,\label{eq:right-raw-transform}
\end{align}
and, in particular,
\begin{equation}\label{eq:D-transform}
 \boxed{D_{A_u(G)}(a_u+ut)=uD_G(t).}
\end{equation}
\end{lemma}

\begin{proof}
By \eqref{eq:norton-recursion}, a Left option of $A_u(G)$ is
$A_u(G^L)+a_u$.  Numeric translation shifts both of its walls by $a_u$.
Therefore \eqref{eq:raw-walls} and the induction hypothesis give
\[
 \max_{G^L}\bigl(RS_{a_u+ut}(A_u(G^L))+a_u\bigr)
 =\max_{G^L}\bigl(uRS_t(G^L)+a_u\bigr)
 =uL_G^{\rm raw}(t)+a_u.
\]
The Right-option calculation is the same with minimum and the shift $-a_u$,
giving \eqref{eq:right-raw-transform}.  Substitute both identities and
$s=a_u+ut$ into the definition \eqref{eq:gap}:
\begin{align*}
 D_{A_u(G)}(s)
 &=\bigl(uL_G^{\rm raw}(t)+a_u\bigr)
   -\bigl(uR_G^{\rm raw}(t)-a_u\bigr)-2(a_u+ut)\\
 &=u\bigl(L_G^{\rm raw}(t)-R_G^{\rm raw}(t)-2t\bigr).
\end{align*}
This is \eqref{eq:D-transform}.
\end{proof}

At positive input temperature, \eqref{eq:D-transform} immediately excludes a
meeting below $a_u$.  Temperature zero is subtler: both sides of
\eqref{eq:D-transform} vanish at $t=0$, so one must rule out a terminal mast
already present below $a_u$.  The following strategy lemma does exactly that.

\begin{lemma}[full-slope corridor]\label{lem:corridor}
Assume $a_u>0$.  Let $X$ be a canonical nonnumber, let $m$ be a number, and
assume that every proper follower $Z$ of $X$ already satisfies the endpoint
identities
\[
 LS_{a_u}(A_u(Z))=uLS(Z),
 \qquad RS_{a_u}(A_u(Z))=uRS(Z).
\]
Suppose that for some $b<a_u$,
\begin{equation}\label{eq:right-corridor}
 RS_s(A_u(X))=um-a_u+s
 \qquad(b\leq s\leq a_u).
\end{equation}
Equivalently, the transformed wall $\mathcal T_X^R(s)$ is the constant $um$
on that interval.  Then Right wins $X-m$ moving first.  The mirror statement
holds when
$LS_s(A_u(X))=um+a_u-s$: Left wins $X-m$ moving first.
\end{lemma}

\begin{proof}
We prove the Right statement by induction on the birthday of $X$; negation
gives the mirror.  Since the wall in \eqref{eq:right-corridor} has the maximal
allowed slope $+1$, it is an unfrozen right scaffold on $[b,a_u)$.
Using \eqref{eq:norton-recursion} there gives
\begin{equation}\label{eq:min-left-walls}
 \min_{Y\in X^R}LS_s(A_u(Y))=um
 \qquad(b\leq s\leq a_u).
\end{equation}
Choose a Right option $Y$ attaining the minimum at $s=b$.  Every left wall is
nonincreasing, while \eqref{eq:min-left-walls} says that each such wall is at
least $um$ throughout the interval.  Hence the chosen wall cannot fall:
\begin{equation}\label{eq:chosen-flat}
 LS_s(A_u(Y))=um\qquad(b\leq s\leq a_u).
\end{equation}

If $Y$ is numeric, Lemma~\ref{lem:numeric-images} says it has already frozen
strictly below $a_u$, and its mast at $a_u$ is $uY$.  Equation
\eqref{eq:chosen-flat} therefore forces $Y=m$.  Right moves from $X-m$ to
$Y-m=0$ and wins.

Suppose $Y$ is nonnumeric.  Let $Z$ be any Left option of $Y$.  Its transformed
right wall is one of the candidates under the flat left wall
\eqref{eq:chosen-flat}; consequently
\begin{equation}\label{eq:reply-bound}
 RS_s(A_u(Z))+a_u-s\leq um
 \qquad(b\leq s\leq a_u).
\end{equation}
At $s=a_u$, the endpoint hypothesis gives
$RS_{a_u}(A_u(Z))=uRS(Z)$.  Thus $RS(Z)\leq m$.

If $RS(Z)<m$, \eqref{eq:stop-implication} says that Right wins $Z-m$ moving
first.  If $RS(Z)=m$, then the right wall's slope bound $+1$ and its endpoint
$um$ give
\[
 RS_s(A_u(Z))\geq um-a_u+s.
\]
Together with \eqref{eq:reply-bound}, this is equality on $[b,a_u]$.
The induction hypothesis applied to the lower-birthday game $Z$ again says
that Right wins $Z-m$ moving first.

Right therefore begins $X-m$ by moving $X$ to $Y$.  After every Left reply
$Y\to Z$, the preceding paragraph supplies Right's winning continuation from
$Z-m$.  Number avoidance says that Left gains nothing by moving in the numeric
component $-m$ instead.  This completes the explicit reply strategy and the
birthday induction.
\end{proof}

\begin{theorem}[exact thermic law]\label{thm:thermic-law}
Let $u>0$ be a dyadic number in canonical form, with $\delta_u$ and $a_u$ as
in \eqref{eq:mesh-shift}.  For every nonnumeric short game $G$ in the finite
thermographic domain,
\begin{equation}\label{eq:main-law}
 \mean(A_u(G))=u\mean(G),
 \qquad
 \temp(A_u(G))=u\temp(G)+a_u.
\end{equation}
More strongly, for every $t\geq0$ the frozen walls obey
\begin{equation}\label{eq:wall-law}
 LS_{a_u+ut}(A_u(G))=uLS_t(G),
 \qquad
 RS_{a_u+ut}(A_u(G))=uRS_t(G).
\end{equation}
For numeric $G=x$, Lemma~\ref{lem:numeric-images} is the exact replacement.
\end{theorem}

\begin{proof}
We prove \eqref{eq:main-law} and \eqref{eq:wall-law} simultaneously by birthday
induction.  Numeric subpositions are supplied by Lemma~\ref{lem:numeric-images};
for nonnumeric options, use the birthday induction.  Thus the hypotheses of
Lemma~\ref{lem:raw-gap} hold for $G$, and the raw-gap identity
\eqref{eq:D-transform} is available.

Put $\tau=\temp(G)$.  If $\tau>0$, then $D_G(0)>0$, hence
\[
 D_{A_u(G)}(a_u)=uD_G(0)>0.
\]
Because the raw gap is nonincreasing, it cannot have closed at any smaller
height.  Equation \eqref{eq:D-transform} then says that its first zero at or
above $a_u$ occurs exactly at $a_u+u\tau$.

If $\tau=0$ and $a_u=0$, the same identity locates the first meeting at zero.
It remains only to treat $\tau=0<a_u$.  Here
$LS(G)=RS(G)=m$ for a number $m$, and \eqref{eq:D-transform} locates a meeting
of the transformed raw scaffolds at $a_u$ with value $um$.  Suppose for a
contradiction that $A_u(G)$ had already frozen at some $b<a_u$.  Its two walls
would then be the constant mast $um$ throughout $[b,a_u]$.

Consider an arbitrary Left option $X$ of $G$.  At $s=a_u$, the induction
hypothesis gives
\[
 RS_{a_u}(A_u(X))=uRS(X),
 \qquad
 \max_{G^L}RS(X)=LS(G)=m.
\]
If $RS(X)<m$, then Right wins $X-m$ moving first by
\eqref{eq:stop-implication}.  If $RS(X)=m$, the candidate
$\mathcal T_X^R(s)$ lies below the frozen parent mast $um$ on $[b,a_u]$ and
equals it at $a_u$.  Since $RS_s(A_u(X))$ has slope at most $+1$,
$\mathcal T_X^R$ is nonincreasing; going backwards from its endpoint shows it
cannot be strictly smaller.  Therefore
\[
 \mathcal T_X^R(s)=um,
 \quad\text{or equivalently}\quad
 RS_s(A_u(X))=um-a_u+s
 \quad(b\leq s\leq a_u).
\]
The endpoint hypotheses of Lemma~\ref{lem:corridor} are exactly the already
proved birthday-induction statements for the proper followers of $X$; the
lemma therefore says that Right wins $X-m$ moving first.

The quantifier is the point: the argument has been applied to \emph{every}
Left option $X=G^L$.  Hence every Left first move in the $G$ component of
$G-m$ loses to Right.  By number avoidance, a move in the number component is
not better.  Right therefore wins $G-m$ playing second, which proves $G\leq m$.
The mirror argument, applied option by option to every Right option $G^R$ and
the left-wall half of Lemma~\ref{lem:corridor}, proves $G\geq m$.  Thus $G=m$,
contrary to the hypothesis that $G$ is nonnumeric.  No premature meeting below
$a_u$ is possible.

We have proved in all cases that the first meeting is
$s=a_u+u\tau$.  At $s=a_u+ut$ before that meeting,
\eqref{eq:left-raw-transform} and \eqref{eq:scaffold-walls} give
\[
 LS_s^{\rm sc}(A_u(G))
 =uL_G^{\rm raw}(t)+a_u-(a_u+ut)
 =uLS_t^{\rm sc}(G),
\]
and the right-wall calculation is identical.  At the meeting, their common
value is $u\mean(G)$; above it both sides freeze to that mast.  This proves
\eqref{eq:wall-law} and both formulas in \eqref{eq:main-law}.
\end{proof}

\begin{corollary}[thermic residue transport]\label{cor:transport}
Let $F_{<0}=\Numbers$ be the subgroup of short numbers and, for $\tau\geq0$,
let $F_{\leq\tau}$ be the subgroup of games of temperature at most $\tau$,
including $F_{<0}$.  Then $A_u$ induces an additive map
\begin{equation}\label{eq:graded-transport}
 \overline A_{u,\tau}:
 F_{\leq\tau}/F_{<\tau}
 \longrightarrow
 F_{\leq u\tau+a_u}/F_{<u\tau+a_u}.
\end{equation}
It is the generalized overheating operator $\int_u^{a_u}$.  For a positive
integer $n$, the shifted thermal height $h=\tau+1$ scales exactly:
$h(A_n(G))=nh(G)$ on nonnumeric layers.
\end{corollary}

\begin{proof}
Norton multiplication by a fixed positive unit is additive
\cite[Prop.~13]{mesdal2009}.  If $G-H$ is nonnumeric of temperature
$\sigma<\tau$, Theorem~\ref{thm:thermic-law} gives
$\temp(A_u(G-H))=u\sigma+a_u<u\tau+a_u$.  If $G-H$ is numeric,
Lemma~\ref{lem:numeric-images} puts its image strictly below $a_u$, hence below
the target layer.  Thus \eqref{eq:graded-transport} is independent of the
representative.  Recurrence \eqref{eq:norton-recursion} is precisely
$\int_u^{a_u}$.  Finally $a_n=n-1$, so
$\temp(A_n(G))+1=n(\temp(G)+1)$.
\end{proof}

\section{The exact composition defect}

Write
\[
 r_u(\tau)=u\tau+a_u
            =u\tau+u-\delta_u.
\]
Corollary~\ref{cor:transport} says that $A_u$ carries thermic degree $\tau$ to
$r_u(\tau)$.  The discrepancy between successive transports and transport by
the ordinary dyadic product is completely arithmetic.

\begin{theorem}[composition defect]\label{thm:defect}
For positive dyadics $u,v$ and every $\tau\geq0$,
\begin{equation}\label{eq:defect}
 r_v(r_u(\tau))=r_{uv}(\tau)+\Delta(u,v),
 \qquad
 \Delta(u,v)=v(1-\delta_u)-\delta_v+\delta_{uv}\geq0.
\end{equation}
The defect vanishes exactly in the following cases:
\begin{enumerate}[label=(\roman*),itemsep=0pt]
\item $u$ is nonintegral and $v=2^{-\ell}$ for some $\ell\geq0$;
\item $u$ is an odd positive integer and $v$ is arbitrary;
\item $u$ is an even positive integer and $v$ is an integer.
\end{enumerate}
\end{theorem}

\begin{proof}
Expanding the affine maps gives
\begin{align*}
 r_v(r_u(\tau))-r_{uv}(\tau)
 &=v(u\tau+u-\delta_u)+v-\delta_v
   -(uv\tau+uv-\delta_{uv})\\
 &=v(1-\delta_u)-\delta_v+\delta_{uv}.
\end{align*}

If $u$ is integral, $\delta_u=1$ and
$\Delta=\delta_{uv}-\delta_v$.  Multiplication by an integer cannot increase
a dyadic denominator.  It preserves the denominator of every $v$ exactly when
$u$ is odd; if $u$ is even and $v$ nonintegral, it cancels at least one factor
of two.  Integral $v$ always has $\delta_{uv}=\delta_v=1$.  This gives (ii),
(iii), and nonnegativity in the integral-$u$ case.

Now write $u=m/2^k$ in lowest terms with $k>0$.  If $v$ is nonintegral, both
numerators are odd and $\delta_{uv}=\delta_u\delta_v$.  Hence
\[
 \Delta=(1-\delta_u)(v-\delta_v)\geq0,
\]
with equality exactly when $v$ has numerator $1$, that is,
$v=2^{-\ell}$.  If $v=n$ is integral, write $n=2^sq$ with $q$ odd.  For
$s\geq k$, one has $\delta_{uv}=1$ and
$\Delta=n(1-2^{-k})>0$ unless $n=1$.  For $s<k$,
\[
 \Delta=n-1-\frac{n-2^s}{2^k}.
\]
This is zero at $n=1$ and positive at every $n>1$: for $s=0$ it is
$(n-1)(1-2^{-k})$, while for $s>0$ the subtracted fraction is at most
$(n-2^s)/2<n-1$.  This proves (i) and completes the classification.
\end{proof}

\begin{corollary}[not a dyadic action]\label{cor:not-action}
No temperature-preserving residue refinement can make all numeric Norton
transports into the scalar action of an associative algebra whose positive
dyadic coefficients multiply ordinarily.
\end{corollary}

\begin{proof}
Take $u=1/2$, $v=2$, and $G=*$.  Theorem~\ref{thm:defect} gives
$\Delta(1/2,2)=1$, hence
\[
 \temp(A_2A_{1/2}(*))=1,
 \qquad
 \temp(A_1(*))=\temp(*)=0.
\]
Associativity would identify the two outputs, but a refinement preserving exact
thermic degree cannot identify classes in different degrees.
\end{proof}

\section{Machine verification in this repository}

The proof above is mathematical; the implementation supplies independent
regression oracles for its formulas.  The following helpers live in
\path{src/games/heating.rs}:
\begin{center}
\small
\path{numeric_norton_regrade}\qquad
\path{numeric_norton_mean_temperature}.
\end{center}
The first returns $(u,a_u)$; the second evaluates both branches of
Theorem~\ref{thm:thermic-law} without materializing the recursive product.
Source tests compare those helpers with actual Norton products on integers,
nonintegral numbers, all-small games, switches, and a nested hot game, using
eight units from $1/4$ through $3$.  A separately transcribed recursive Norton
oracle checks the noninteger-game/noninteger-unit branch.

The maintained probe \path{experiments/under_descent.py} uses a $35$-game
catalogue and six units.  It checks $210$ exact mean/temperature predictions,
$24$ representative-pair descent tests for Norton multiplication, and $24$ for
the matching overheating operator.  Its arithmetic defect scan checks $2{,}304$
positive-dyadic pairs for nonnegativity and the complete zero classification,
then compares five selected defects with recursively materialized game
products.  It also records the explicit nonnumeric-unit failures discussed
below.  These are bounded machine checks, not substitutes for the universal
proofs.

\section{The tropical shadow and its limits}

\subsection{What is genuinely tropical}

Temperature satisfies the standard maximum inequality
\begin{equation}\label{eq:max-inequality}
 \temp(G+H)\leq\max(\temp(G),\temp(H)),
\end{equation}
with numbers colder than every infinitesimal or hot game
\cite[Thm.~II.5.18]{siegel2013}; see also the cooling-homomorphism treatment in
\cite{bando2020}.  Consequently the $F_{\leq\tau}$ above are additive
subgroups and define
\[
 \gr_T(\Games)=\bigoplus_{\tau\geq0}F_{\leq\tau}/F_{<\tau}.
\]
Equal leading temperatures may cancel to a lower layer, just as equal
valuations may cancel in a valued field.

There is also an exact max-plus/min-plus observation inside the wall recursion:
$L_G^{\rm raw}$ is the pointwise maximum of the Left options' right walls, and
$R_G^{\rm raw}$ is the pointwise minimum of the Right options' left walls.
Cooling shifts wall values by ordinary addition.  These are dual tropical
folds.  The subsequent operation is not tropical, however: one subtracts the
two envelopes and $2t$, finds the first zero of $D_G$, and freezes both walls
to a common mast.  The module
\path{src/games/tropical_thermography.rs} deliberately reuses this ordinary
freeze tail; it only names the two option folds.  Thus thermography contains
tropical operations but is not, as a whole, a semiring tropicalization.

The valuation-theoretic comparison remains useful.  A Newton polygon takes the
lower convex hull of $(i,v(a_i))$, and Dumas's theorem makes the polygon of a
product the appropriate sum of the factor polygons \cite{dumas1906}.  The game
side has the filtration inequality \eqref{eq:max-inequality} and the external
regradings \eqref{eq:graded-transport}, but no corresponding internal product
on Conway game values.

Two terminological distinctions prevent predictable misreadings.  First,
``tropical games'' often means the algorithmic equivalence between tropical
polyhedra and finite mean-payoff games established by
Akian--Gaubert--Guterman \cite{akian2012}.  Those are weighted graph
optimization games, not Conway values or thermographs.  Second, ``temperature''
in idempotent analysis commonly denotes the Maslov dequantization parameter
whose zero limit turns log-sum-exp into max-plus arithmetic
\cite{viro2010}.  Combinatorial-game temperature instead measures the urgency
at which cooling freezes a game.  The shared words do not identify the two
constructions.

\subsection{Thermographs and nonnumeric units do not descend}

The exact thermograph is not a congruence for disjunctive sum.  Indeed,
$*=\{0\mid0\}$ and $\uparrow=\{0\mid*\}$ have the same constant-zero walls,
mast zero, and temperature zero.  Yet $*+*=0$ has temperature $-1$, whereas
$\uparrow+*$ is a nonnumber of temperature zero.  Hence no binary operation on
thermographs alone can recover the thermograph of every sum.

Nor does arbitrary Norton multiplication descend to $\gr_T$.  At degree zero,
$*$ and $*+1$ differ by the cold number $1$.  For the positive nonnumeric unit
$U=\uparrow$, additivity and the integer clause give the two-line obstruction
\[
 A_U(*+1)-A_U(*)=A_U(1)=U=\uparrow.
\]
The output difference has a nonzero temperature-zero class.  Thus numericity
of $u$ in Corollary~\ref{cor:transport} is a real boundary, not a proof
convenience.

\subsection{The degree-zero torsion remark}

The degree-zero group is already familiar.  Every temperature-zero game is a
number plus an infinitesimal, and quotienting by the cold numbers gives
\begin{equation}\label{eq:gr0-inf}
 \gr_0(\Games)=F_{\leq0}/F_{<0}\cong\Inf,
\end{equation}
the additive group of infinitesimals.  Calistrate's decomposition
$\Games=\Inf\oplus\mathrm{Rcf}$ places this group inside the classical
structure theory \cite{calistrate1996}.  Its torsion is not created by the
temperature filtration: $[*]\neq0$ and $2[*]=0$ reflect the pre-existing
additive group of games, studied abstractly by Moews \cite{moews2002}.
Once \eqref{eq:gr0-inf} is stated, the coefficient obstruction is the elementary
fact that inverting $2$ kills nonzero $2$-torsion.  We record the precise
contract only to prevent a grading loophole.

\begin{definition}[full-dyadic lift-compatible action]\label{def:full-dyadic}
Such an action consists of the following data.
\begin{enumerate}[label=(\roman*),itemsep=1pt]
\item An abelian grading group $\Gamma$, a degree map
      $\iota_T:\Q_{\geq0}\to\Gamma$, and a $\Gamma$-graded additive group
      $R=\bigoplus_{\gamma\in\Gamma}R_\gamma$.
\item Injective homomorphisms
      $j_\tau:F_{\leq\tau}/F_{<\tau}\hookrightarrow
      R_{\iota_T(\tau)}$ for every thermic degree $\tau$.
\item A graded unital ring $C=\bigoplus_d C_d$ acting associatively and
      distributively on $R$, with
      $C_dR_\gamma\subseteq R_{d+\gamma}$.
\item Mutually inverse homogeneous elements
      $\pi_2\in C_d$ and $\pi_{1/2}\in C_{-d}$ representing the full dyadic
      coefficients $2$ and $1/2$.
\item A common $\Gamma$-filtered lift of the game group extending the thermic
      filtration, such that $j_\tau([G]_\tau)$ is the
      $\iota_T(\tau)$-initial class of $G$ and, in the single component
      $R_{\iota_T(\tau)+d}$,
      \begin{equation}\label{eq:compatibility}
       \pi_2j_\tau([G]_\tau)
       =\operatorname{in}_{\iota_T(\tau)+d}(G+G).
      \end{equation}
      The right side is zero when $G+G$ lies strictly below that common target
      degree, and in particular when $G+G=0$.
\end{enumerate}
The common group $\Gamma$ and the component declaration in
\eqref{eq:compatibility} are part of the data; neither side may be silently
regraded into a different group.
\end{definition}

\begin{proposition}[dyadic torsion no-go]\label{prop:torsion}
No faithful full-dyadic lift-compatible action exists.
\end{proposition}

\begin{proof}
Let $s=j_0([*])\neq0$.  Since $*+*=0$, compatibility gives $\pi_2s=0$.
But $\pi_2$ is invertible, so
\[
 s=(\pi_{1/2}\pi_2)s=\pi_{1/2}(\pi_2s)=0,
\]
a contradiction.
\end{proof}

This does not forbid three alternatives: integer or valuation-ring coefficients
in which $2$ is not inverted; a characteristic-$2$ quotient that kills the
incompatible coefficient data; or a ring built from a finer equivalence
relation rather than Conway equality.  The recent construction of
Altman--Lipparini belongs to the last category \cite{altman2026}.  It is
adjacent to, but does not contradict, Proposition~\ref{prop:torsion}.

\section{Conclusion}

The positive theorem is the durable bridge.  A canonical positive dyadic unit
$u$ sends every nonnumeric thermic degree by the exact affine rule
$\tau\mapsto u\tau+u-\delta_u$, while numeric inputs occupy a strictly lower
branch.  This yields honest additive maps between thermic residue groups.  The
composition defect then gives the exact limit: these maps form a family of
external regradings, not an action of the multiplicative dyadics.

The Newton-polygon comparison should be kept at that scale.  Temperature obeys
a nonarchimedean maximum inequality; the raw thermograph recursion contains
dual tropical folds.  The freeze operation is not tropical, the complete
thermograph is not a sum congruence, and the degree-zero residual group retains
the old infinitesimal torsion.  The game and place sides therefore share a
filtered shadow, not one faithful full-dyadic graded ring.

% REF: berlekamp2001 = Elwyn R. Berlekamp, John H. Conway, Richard K. Guy, Winning Ways for Your Mathematical Plays, 2nd ed., 4 vols., A K Peters, 2001-2004.
% REF: siegel2013 = Aaron N. Siegel, Combinatorial Game Theory, Graduate Studies in Mathematics 146, American Mathematical Society, 2013.
% REF: johnson2011 = Will Johnson, Combinatorial Game Theory, Well-Tempered Scoring Games, and a Knot Game, arXiv:1107.5092, 2011.
% REF: mesdal2009 = G. A. Mesdal III, Partizan Splittles, Games of No Chance 3, MSRI Publications 56, Cambridge University Press, 2009, 447-462.
% REF: bando2020 = Katsuyuki Bando, Eitetsu Ken, Kota Morikawa, Foundations of Temperature Theory, arXiv:2009.02876, 2020.
% REF: calistrate1996 = Dan Calistrate, The Reduced Canonical Form of a Game, Games of No Chance, MSRI Publications 29, Cambridge University Press, 1996, 409-416.
% REF: moews2002 = David Moews, The Abstract Structure of the Group of Games, More Games of No Chance, MSRI Publications 42, Cambridge University Press, 2002, 49-57.
% REF: dumas1906 = Gustave Dumas, Sur quelques cas d'irreductibilite des polynomes a coefficients rationnels, Journal de Mathematiques Pures et Appliquees, serie 6, 2 (1906), 191-258.
% REF: akian2012 = Marianne Akian, Stephane Gaubert, Alexander Guterman, Tropical polyhedra are equivalent to mean payoff games, International Journal of Algebra and Computation 22:1 (2012), 1250001.
% REF: viro2010 = Oleg Viro, Hyperfields for tropical geometry I: Hyperfields and dequantization, arXiv:1006.3034, 2010.
% REF: altman2026 = Harry Altman, Paolo Lipparini, A Ring structure on the Class of Combinatorial Games, arXiv:2604.27847, 2026.

\footnotesize
\begin{thebibliography}{99}

\bibitem{akian2012}
Marianne Akian, St\'ephane Gaubert, and Alexander Guterman,
``Tropical polyhedra are equivalent to mean payoff games,''
\emph{International Journal of Algebra and Computation} 22:1 (2012), 1250001.

\bibitem{altman2026}
Harry Altman and Paolo Lipparini,
``A Ring structure on the Class of Combinatorial Games,''
arXiv:2604.27847 (2026).

\bibitem{bando2020}
Katsuyuki Bando, Eitetsu Ken, and Kota Morikawa,
``Foundations of Temperature Theory,'' arXiv:2009.02876 (2020).

\bibitem{berlekamp2001}
Elwyn R. Berlekamp, John H. Conway, and Richard K. Guy,
\emph{Winning Ways for Your Mathematical Plays}, 2nd ed., 4 vols.,
A K Peters, 2001--2004.

\bibitem{calistrate1996}
Dan Calistrate,
``The Reduced Canonical Form of a Game,'' in
\emph{Games of No Chance}, MSRI Publications 29,
Cambridge University Press, 1996, 409--416.

\bibitem{dumas1906}
Gustave Dumas,
``Sur quelques cas d'irr\'eductibilit\'e des polyn\^omes \`a coefficients
rationnels,'' \emph{Journal de Math\'ematiques Pures et Appliqu\'ees},
6e s\'erie, 2 (1906), 191--258.

\bibitem{johnson2011}
Will Johnson,
\emph{Combinatorial Game Theory, Well-Tempered Scoring Games, and a Knot Game},
arXiv:1107.5092 (2011).

\bibitem{mesdal2009}
G. A. Mesdal III,
``Partizan Splittles,'' in \emph{Games of No Chance 3},
MSRI Publications 56, Cambridge University Press, 2009, 447--462.

\bibitem{moews2002}
David Moews,
``The Abstract Structure of the Group of Games,'' in
\emph{More Games of No Chance}, MSRI Publications 42,
Cambridge University Press, 2002, 49--57.

\bibitem{siegel2013}
Aaron N. Siegel,
\emph{Combinatorial Game Theory}, Graduate Studies in Mathematics 146,
American Mathematical Society, 2013.

\bibitem{viro2010}
Oleg Viro,
``Hyperfields for tropical geometry I: Hyperfields and dequantization,''
arXiv:1006.3034 (2010).

\end{thebibliography}

\end{document}
